% Informe técnico del Reto III.3. Compilar con XeLaTeX o LuaLaTeX (Overleaf: Menú > Compilador).
% Las figuras están en la carpeta figuras/ (se generan con: python reproducir.py --reporte).
\documentclass[11pt,letterpaper]{article}

\usepackage{iftex}
\ifPDFTeX
  \usepackage[utf8]{inputenc}
  \usepackage[T1]{fontenc}
  \usepackage{lmodern}
  \usepackage[spanish]{babel}
\else
  \usepackage{fontspec}
  \usepackage{polyglossia}
  \setmainlanguage{spanish}
\fi

\usepackage[margin=2.3cm,headheight=15pt]{geometry}
\usepackage{amsmath,amssymb}
\usepackage{graphicx}
\usepackage{booktabs}
\usepackage{array}
\usepackage{siunitx}
\usepackage[font=small,labelfont=bf]{caption}
\usepackage{fancyhdr}
\usepackage{xcolor}
\usepackage[hidelinks]{hyperref}
\setlength{\emergencystretch}{3em}

\newcommand{\RJ}{R_{\mathrm J}}
\newcommand{\vinf}{v_\infty}
\newcommand{\rsoi}{r_{\mathrm{SOI}}}

\definecolor{enlace}{RGB}{0,70,160}
\newcommand{\vlink}[2]{\href{https://drive.google.com/file/d/#1/view?usp=sharing}{\textcolor{enlace}{\underline{#2}}}}
\newcommand{\video}[2]{\par\nopagebreak\smallskip\noindent\hfill{\small $\triangleright$ Video de referencia: \vlink{#1}{#2}}\par\smallskip}
\pagestyle{fancy}
\fancyhf{}
\fancyhead[C]{Asistencia gravitacional en Júpiter}
\fancyhead[R]{\small noviembre de 2026}
\fancyfoot[C]{\small Universidad Industrial de Santander}
\fancyfoot[R]{\thepage}
\renewcommand{\footrulewidth}{0.4pt}

\title{Diseño de una asistencia gravitacional en Júpiter:\\
cónicas empalmadas frente a integración de $N$ cuerpos}
\author{Frankly Giovanni Poveda Pinzón
\quad\quad Camila Castillo López\\[4pt]
Universidad Industrial de Santander\\
Escuela de Física, Bucaramanga, Colombia}
\date{18 de noviembre de 2026}

\begin{document}
\renewcommand{\contentsname}{Índice}
\maketitle
\thispagestyle{fancy}



\setcounter{tocdepth}{1}
\tableofcontents

\begin{abstract}
Se diseñó un sobrevuelo de Júpiter hacia Saturno con dos modelos: cónicas empalmadas e
integración de $N$ cuerpos. A partir del lagrangiano en el sistema que rota con Júpiter se
muestra que la energía y el momento angular de la sonda respecto al Sol no se conservan por
separado, pero sí $h = E - n\ell_z$, de donde $\Delta E = n\,\Delta\ell_z$; en el sistema de
Júpiter se conservan, aproximadamente, su energía y su momento angular, y con ellos $\vinf$. Con
efemérides reales JPL DE421, las cónicas encuentran la ventana de 1977 y predicen el periapsis
de las Voyager con error menor al $4\,\%$. Su error en la ganancia de energía es menor al
$2\,\%$ para $\vinf \gtrsim 8$~km/s y supera el $10\,\%$ para $\vinf \lesssim 3$~km/s. Al
integrar Voyager~2 desde su estado real, el modelo de $N$ cuerpos solo reproduce la trayectoria
si se incluyen los satélites galileanos, el achatamiento de Júpiter y una maniobra de la sonda
de 76 minutos, detectada con la aceleración residual.
\end{abstract}



\section{Introducción}

Para llegar a los planetas exteriores se necesita mucha energía. Una transferencia de Hohmann
directa a Saturno tarda unos seis años. La alternativa, propuesta en los años sesenta
\cite{flandro1966}, es el sobrevuelo planetario o asistencia gravitacional: la sonda pasa cerca
de un planeta y sale con otra velocidad respecto al Sol sin gastar combustible. En 1977 una
alineación favorable permitió que las dos sondas Voyager usaran a Júpiter para llegar a
Saturno en tres a cuatro años \cite{kohlhase1977}.

El fenómeno parece contradecir la conservación de la energía, porque la sonda sale de Júpiter
más rápido de lo que llegó. La contradicción se resuelve al especificar el sistema de
referencia, y esa es la primera parte de la pregunta del reto: \emph{qué cantidad debe
conservarse durante el encuentro en cada sistema de referencia}. La segunda parte es
cuantitativa: \emph{cuál es el error de la aproximación de cónicas empalmadas}, el método que
divide la trayectoria en tramos donde actúa un solo cuerpo y que se usa en el diseño
preliminar de misiones \cite{curtis2020}.

El caso de estudio es el de las Voyager: Voyager~2 pasó a unos $10\,\RJ$ de Júpiter en julio de
1979 y llegó a Saturno en 1981. Ese encuentro está documentado con efemérides precisas, de modo
que sirve para medir cuánto se equivoca cada modelo. Se comparan las cónicas empalmadas, que son
analíticas pero suponen un solo cuerpo en cada tramo, con la integración numérica de los $N$
cuerpos, que no hace esa suposición.

Este trabajo responde ambas preguntas con mecánica clásica: la sección~\ref{sec:metodo} presenta la formulación, los modelos, los datos y las hipótesis; la~\ref{sec:resultados}, los resultados y sus verificaciones; la~\ref{sec:conclusiones}, lo que es confiable.

\section{Metodología}\label{sec:metodo}

\subsection{Formulación lagrangiana y cantidades conservadas}

Se considera una sonda de masa despreciable atraída por el Sol ($M_1$) y por Júpiter ($M_2$),
que giran en órbitas circulares alrededor de su centro de masa con velocidad angular $n$,
$n^2a^3 = G(M_1+M_2)$. En un sistema que rota con Júpiter, con el Sol en $(-\alpha a,0,0)$ y
Júpiter en $((1-\alpha)a,0,0)$, $\alpha = M_2/(M_1+M_2)$, el lagrangiano por unidad de masa es
\begin{equation}
L = \tfrac12\left[(\dot x - ny)^2 + (\dot y + nx)^2 + \dot z^2\right] + \frac{GM_1}{r_1} + \frac{GM_2}{r_2},
\end{equation}
donde los paréntesis son las componentes de la velocidad inercial. Las ecuaciones de
Euler--Lagrange dan $\ddot x - 2n\dot y = \partial_x\Omega$, $\ddot y + 2n\dot x = \partial_y\Omega$,
$\ddot z = \partial_z\Omega$, con $\Omega = \tfrac12 n^2(x^2+y^2) + GM_1/r_1 + GM_2/r_2$; los
términos $\pm 2n\dot q$ son la fuerza de Coriolis. Como $L$ no depende explícitamente del
tiempo, se conserva la función de energía \cite{goldstein2002}
\begin{equation}
h = \sum_k \dot q_k\frac{\partial L}{\partial\dot q_k} - L
  = \tfrac12\left(\dot x^2+\dot y^2+\dot z^2\right) - \Omega = E - n\,\ell_z ,
\label{eq:jacobi}
\end{equation}
donde $E$ y $\ell_z$ son la energía y el momento angular de la sonda medidos en el sistema
inercial; $C_J = -2h$ es la constante de Jacobi \cite{murray1999}. En el sistema inercial el
potencial depende del tiempo y no tiene simetría de rotación, así que ni $E$ ni $\ell_z$ se
conservan; solo lo hace la combinación~\eqref{eq:jacobi}, asociada a la simetría ante una
traslación temporal acompañada de una rotación $n\Delta t$. De ahí,
\begin{equation}
\Delta E = n\,\Delta\ell_z .
\label{eq:DE}
\end{equation}
Esta igualdad exige que la órbita de Júpiter sea circular. Si es excéntrica, $h$ no es
constante y $dE/dt = \omega(t)\,d\ell_z/dt$ más un término proporcional a la velocidad radial de
Júpiter, donde $\omega(t)$ es su velocidad angular instantánea. Lejos de Júpiter,
\eqref{eq:jacobi} se reduce al parámetro de Tisserand
$T_{\mathrm J} = a/a_s + 2\cos i\sqrt{(a_s/a)(1-e^2)}$, donde $a_s$, $e$, $i$ son los elementos
heliocéntricos de la sonda.

En el sistema de Júpiter, la ecuación de la posición relativa $\vec\rho$ contiene la atracción
de Júpiter más una perturbación solar de marea de orden $GM_1\rho/a^3$. Si esta se desprecia, el
problema es de dos cuerpos: se conservan la energía y el momento angular jovicéntricos, la
órbita es una hipérbola, $|\vec v_{\infty,\mathrm{ent}}| = |\vec v_{\infty,\mathrm{sal}}|$ y la
velocidad gira un ángulo $\delta$ dado por
\begin{equation}
e = 1 + \frac{r_p\vinf^2}{GM_2}, \qquad \sin\frac{\delta}{2} = \frac1e .
\label{eq:giro}
\end{equation}
Igualando el cociente perturbación/atracción en ambos sistemas se obtiene el radio de la
esfera de influencia de Laplace, $\rsoi = a(M_2/M_1)^{2/5} = 674\,\RJ$. La deducción completa
está en el documento de formulación del repositorio, verificada con álgebra simbólica (sympy).

\subsection{Modelo 1: cónicas empalmadas}

Los tramos Tierra--Júpiter y Júpiter--Saturno se obtienen resolviendo el problema de Lambert
con variable universal \cite{curtis2020}; el encuentro se trata como instantáneo visto desde el
Sol. Para cada fecha de salida y tiempos de vuelo se calcula el $\vec v_\infty$ de llegada y el
que exige el segundo tramo. El sobrevuelo es posible sin motor si sus magnitudes coinciden y si
el periapsis que da~\eqref{eq:giro} supera $2\,\RJ$. Se barrieron $4.2\times10^6$ combinaciones
entre 1976 y 1980 y las mejores se refinaron con optimización con restricciones (SLSQP),
minimizando el tiempo de viaje para un límite de la energía de lanzamiento $C_3$.

\subsection{Modelo 2: integración de $N$ cuerpos}

Se integraron las ecuaciones de Newton con REBOUND y el integrador de paso adaptativo IAS15
\cite{rein2012,rein2015}. Para medir el error de las cónicas se usó el problema restringido
circular (Sol, Júpiter en órbita circular y sonda), en el cual \eqref{eq:jacobi} es exacta y
sirve de control. Ambos modelos parten del \emph{mismo estado en el periapsis} y se propagan
hacia atrás y hacia adelante hasta que la sonda está a más de $3\,\rsoi$ de Júpiter. En el
Modelo 1 la sonda sigue la hipérbola hasta el radio de empalme $R_e$ y después una cónica
heliocéntrica. Como medida del error se tomó la diferencia relativa en la ganancia de energía
heliocéntrica $\Delta E$, que es la cantidad física que produce el sobrevuelo. Se barrieron
$\vinf$ entre 2.5 y 15~km/s, $R_e$ entre $0.25$ y $3\,\rsoi$ y $r_p = 3$, 10 y $30\,\RJ$.

Para comparar con Voyager~2 se usaron dos versiones. La básica tiene al Sol, a Júpiter como
masa puntual y a Saturno. En la ampliada Júpiter se separa en el planeta, con su achatamiento
$J_2$, y los cuatro satélites galileanos; las posiciones de todos los cuerpos se toman de las
efemérides y solo se integra la sonda (DOP853, tolerancia relativa $10^{-12}$).

\subsection{Datos reales}

Todos los datos son públicos y no se ajustaron en este trabajo. El cuadro~\ref{tab:datos} dice de dónde sale cada uno; \texttt{datos/LEEME.md} del repositorio da el procedimiento y las comprobaciones. El kernel de Voyager~2 se validó con valores publicados \cite{pds}: máximo acercamiento a Júpiter el 9 de julio de 1979 a las 22:29:01~UTC, a 721\,673~km (publicado: 721\,670~km), y distancias a los cuatro satélites galileanos al $0.2\,\%$ o mejor (Ganímedes: 62\,078 frente a 62\,130~km).

\begin{table}[ht]\centering\footnotesize\renewcommand{\arraystretch}{1.05}
\caption{Origen de los datos.}\label{tab:datos}
\begin{tabular}{>{\raggedright\arraybackslash}p{4.1cm}>{\raggedright\arraybackslash}p{7.3cm}>{\raggedright\arraybackslash}p{3.1cm}}
\toprule
Dato & Fuente & Uso\\\midrule
Efemérides de Sol, Tierra, Júpiter y Saturno; $GM$ de Sol, Júpiter y Saturno & JPL DE421 \cite{folkner2009} (paquete \texttt{de421}, leído con \texttt{jplephem}) & Ambos modelos\\
Radio (71\,492~km) y polo de Júpiter & Unión Astronómica Internacional \cite{archinal2018} & $r_p$ en $R_J$; $J_2$\\
$GM$ de los satélites galileanos y $J_2$ de Júpiter & Solución JUP230 del JPL \cite{jacobson} & Modelo 2 ampliado\\
Posiciones de los satélites galileanos & Astronomy Engine 2.1.19 \cite{cross} & Modelo 2 ampliado\\
Trayectoria de Voyager~2, 4 jun--20 jul 1979 & Kernel SPICE \texttt{vgr2\_jup230.bsp} de NAIF \cite{acton1996}, copia idéntica del paquete \texttt{asteRisk} & Solo validación\\
Fechas y distancias publicadas & PDS Ring-Moon Systems Node \cite{pds} & Comparación\\
Maniobra de Voyager~2 (76 min) & NASA SP-439, cap.~7 \cite{morrison1980} & Contraste\\
Órbitas circulares (alternativa) & Elementos medios de Standish \cite{standish} & Cuadros 1 y 2\\
\bottomrule
\end{tabular}
\end{table}

\subsection{Hipótesis y verificación}

El cuadro~\ref{tab:hip} resume las decisiones de modelado más importantes; la hoja completa,
con doce hipótesis y cinco decisiones numéricas, está en el repositorio. Antes de simular se registraron tres predicciones:
el giro máximo con $r_p \geq 2\,\RJ$ y $\vinf\approx 8$~km/s sería de unos \ang{138}, la sonda
ganaría energía si pasa por detrás de Júpiter, y el error de las cónicas crecería al bajar
$\vinf$. Cada resultado se verificó por al menos un mecanismo independiente: un segundo integrador (DOP853), invariantes, límites conocidos ($M_2\to0$, $r_p\to\infty$), análisis dimensional y datos no ajustados.

\begin{table}[tb]
\centering
\small
\caption{Hipótesis principales y su efecto medido sobre el resultado.}
\label{tab:hip}
\begin{tabular}{>{\raggedright}p{2.6cm}>{\raggedright}p{3.2cm}>{\raggedright}p{2.7cm}>{\raggedright\arraybackslash}p{6.3cm}}
\toprule
Decisión & Elegido & Alternativa & Efecto medido \\
\midrule
Órbitas planetarias & Efemérides DE421 & Circulares & Error en $r_p$ de Voyager~2: $25\,\%$ (circular) frente a $3.7\,\%$ (DE421) \\
Radio de empalme & Laplace, $674\,\RJ$ & $0.25$--$3\,\rsoi$, Hill & Con $\vinf=5$ km/s: error de 27\,\% ($0.25\,\rsoi$) a 2.4\,\% ($2\,\rsoi$); 5.7\,\% con Laplace \\
Duración del encuentro & Instantáneo & Integrado & Error de 17.5\,\% a 0.13\,\% al pasar de $\vinf=2.5$ a 15 km/s \\
Júpiter & Masa puntual & Planeta con $J_2$ y satélites & Voyager~2 real: error en el periapsis de 1300 km (puntual) frente a 310 km; los satélites cambian $\vinf$ en $-0.03$ km/s \\
Saturno en el encuentro & Despreciado & Incluido & Cambia $\Delta E$ en $1$--$4\times10^{-4}$ relativo \\
Sobrevuelo & Sin motor & Con maniobra & Define las fechas factibles. Voyager~2 hizo una de 13 m/s ($0.2\,\%$ de $\vinf$) \\
\bottomrule
\end{tabular}
\end{table}

\section{Resultados y verificación}\label{sec:resultados}

\subsection{Búsqueda de trayectorias y comparación con las Voyager}

Con efemérides reales, las soluciones aparecen en ventanas estrechas, una por año
(figura~\ref{fig:ventanas}a). Las fechas reales de ambas Voyager caen dentro de la ventana de
1977 que el modelo encuentra sin haberlas usado (figura~\ref{fig:ventanas}b). Con
$C_3\leq100$~km$^2$/s$^2$ el viaje más corto sale el 3 de septiembre de 1977, pasa por Júpiter
el 12 de abril de 1979 con $\vinf = 9.84$~km/s y llega a Saturno el 1 de febrero de 1981
(3.42 años).

\video{1a6Vb9RsPnKq5cdguYI7ysPRqMSLKdguq}{la misión Tierra--Saturno}
\begin{figure}[tb]
\centering
\includegraphics[width=0.78\textwidth]{figuras/fig01_ventanas.png}
\caption{Ventanas Tierra--Júpiter--Saturno con sobrevuelo sin motor, efemérides DE421. Color:
tiempo mínimo hasta Saturno con $C_3 \leq 105$~km$^2$/s$^2$; blanco: sin solución; líneas
grises: $C_3$ de salida. (a) 1976--1980. (b) Ventana de 1977 con las fechas reales de Voyager~2
(círculo) y Voyager~1 (cuadrado).}
\label{fig:ventanas}
\end{figure}

El cuadro~\ref{tab:voy} compara lo que el modelo exige en las fechas reales con lo que
hicieron las sondas. La hipótesis de órbitas circulares es la de mayor impacto: sin ella, el
periapsis predicho queda a menos del $4\,\%$ del real. El giro predicho para Voyager~2,
\ang{96.8}, coincide con el obtenido de los datos (\ang{96.5} entre los extremos del registro,
\ang{97.6} con la hipérbola osculante en el periapsis), y el desajuste entre las $\vinf$ de
entrada y salida es de solo $0.10$~km/s: en esas fechas el sobrevuelo se puede hacer
prácticamente sin propulsión.

\begin{table}[tb]
\centering
\small
\caption{Modelo 1 evaluado en las fechas reales de las Voyager. $r_p$ es la distancia de
máximo acercamiento requerida para el sobrevuelo sin motor.}
\label{tab:voy}
\begin{tabular}{lccccc}
\toprule
& $r_p$ circular & $r_p$ DE421 & $r_p$ real & $\delta$ DE421 & $\vinf$ entrada DE421 \\
\midrule
Voyager 2 & $12.59\,\RJ$ & $9.72\,\RJ$ & $10.09\,\RJ$ & \ang{96.8} & 7.89 km/s \\
Voyager 1 & $5.75\,\RJ$ & $4.74\,\RJ$ & $4.88\,\RJ$ & \ang{98.4} & 10.95 km/s \\
\bottomrule
\end{tabular}
\end{table}

\subsection{Qué se conserva en cada sistema}

La figura~\ref{fig:sobrevuelo} muestra el encuentro óptimo en los dos sistemas. En el de
Júpiter, la sonda describe una hipérbola y sale con la misma $\vinf$ con la que entró. En el del
Sol, el giro de $\vec v_\infty$ hacia la dirección de movimiento de Júpiter aumenta la rapidez
de 11.6 a 22.2~km/s: la energía heliocéntrica no se conserva.

\video{1A9H1X37hOCKHA-NX3c8qpK5j5CFkZ6PT}{el encuentro visto desde Júpiter y desde el Sol}
\begin{figure}[tb]
\centering
\includegraphics[width=0.70\textwidth]{figuras/fig04_sobrevuelo.png}
\caption{Sobrevuelo óptimo con $C_3\leq100$~km$^2$/s$^2$. (a) Sistema de Júpiter: hipérbola con
$r_p = 5.96\,\RJ$ y $\delta = \ang{97.9}$. (b) Sistema del Sol: $\vec v = \vec v_{\mathrm J} +
\vec v_\infty$; la velocidad de salida termina sobre el círculo de radio $\vinf$ porque
$|\vec v_\infty|$ se conserva.}
\label{fig:sobrevuelo}
\end{figure}

La cantidad que sí se conserva es $h$. En el Modelo 2, durante un sobrevuelo de 500 días en el
que la energía heliocéntrica pasa de $-157$ a $-40$~km$^2$/s$^2$, $h$ varía en
$2.5\times10^{-14}$ relativo. Con los datos reales de Voyager~2, entre el 4 de junio y el 20 de
julio de 1979, la energía de la sonda aumenta $\Delta E = 164.2$~km$^2$/s$^2$. Con el movimiento
medio de Júpiter, $n\Delta\ell_z = 173.2$~km$^2$/s$^2$: la relación~\eqref{eq:DE} falla un
$5.5\,\%$. La causa es la excentricidad de Júpiter ($0.048$): en julio de 1979 estaba a 5.33~UA
del Sol, cerca del afelio, y su velocidad angular era $0.952\,n$. Con la velocidad angular
instantánea, $\int\omega\,d\ell_z = 164.8$~km$^2$/s$^2$ y la diferencia baja al $0.4\,\%$; el
resto lo explican la velocidad radial de Júpiter y las fuerzas que no están en el potencial
(sección~\ref{sec:voy}). Por la misma razón el
parámetro de Tisserand, que se deduce de $h$ con el $n$ medio, tampoco se conserva: prolongando
la trayectoria con el Modelo 2 hasta más de $3\,\rsoi$ a cada lado, pasa de 2.576 a 2.670
($3.6\,\%$). Lejos de Júpiter la energía heliocéntrica pasa de $-121.8$ a $+26.3$~km$^2$/s$^2$: la
sonda queda en una órbita de escape del Sistema Solar, con excentricidad 1.30.

\video{1aP_AyXtxezNrNl6CwevwEhPgkhLemugn}{$E$, $n\ell_z$ y $h$ durante el sobrevuelo}
\subsection{Verificación del Modelo 2}

Al cambiar la tolerancia de IAS15 entre $10^{-4}$ y $10^{-11}$, la posición final tras 250 días
cambia como mucho 2.1~m (con $\epsilon = 10^{-4}$; menos de 2~cm con $\epsilon \le 10^{-6}$), y un integrador independiente (DOP853 de SciPy) difiere en 3~m. Cerca del
periapsis ambos modelos coinciden (límite en que domina Júpiter), y con la masa de Júpiter
tendiendo a cero la energía heliocéntrica no cambia. Las 40 pruebas automáticas del repositorio
cubren estos casos.

\subsection{Error de las cónicas empalmadas}

La figura~\ref{fig:error} es el resultado central. Con empalme en la esfera de Laplace y
$r_p = 10\,\RJ$, el error en la ganancia de energía es de $17.5\,\%$ para $\vinf = 2.5$~km/s,
$5.7\,\%$ para 5~km/s, $2.0\,\%$ para 8~km/s, $1.0\,\%$ para 10~km/s y $0.13\,\%$ para 15~km/s, y las
cónicas siempre sobreestiman la ganancia. Esto confirma la tercera predicción: con $\vinf$
pequeña la sonda pasa más tiempo cerca del borde de la esfera, donde ni el Sol ni Júpiter son
despreciables. El error también crece al disminuir $r_p$.

\video{1RBG6jed7L2yAaoZgWtrvhcYv0UIYJz84}{el error de las cónicas según $\vinf$}
\begin{figure}[tb]
\centering
\includegraphics[width=0.85\textwidth]{figuras/fig07_error_conicas.png}
\caption{Error relativo de las cónicas en la ganancia de energía heliocéntrica frente al
modelo de $N$ cuerpos, partiendo del mismo periapsis. (a) Según el radio de empalme, con
$r_p=10\,\RJ$. (b) Según el periapsis, con empalme en $\rsoi$. Puntos llenos: las cónicas
sobreestiman la ganancia; vacíos: la subestiman.}
\label{fig:error}
\end{figure}

La elección del radio de empalme es un compromiso. Con radios pequeños el modelo ignora a
Júpiter fuera de la esfera y el error llega al $70\,\%$; con radios grandes ignora al Sol dentro
de ella, el error cambia de signo para $\vinf$ entre 6 y 7~km/s y la violación de la constante de
Jacobi por las cónicas crece: hasta $0.07\,\%$ con $0.25\,\rsoi$, $0.4\,\%$ con Laplace y $3\,\%$
con $3\,\rsoi$. Laplace y Hill quedan cerca del mejor compromiso.

\subsection{Voyager 2 con ambos modelos}\label{sec:voy}

Partiendo del estado real del 4 de junio de 1979 (a $0.58\,\rsoi$ de Júpiter), el Modelo 2
básico sigue la trayectoria real con 1300~km de error en el periapsis, mientras que la hipérbola
kepleriana se desvía 100\,000~km. Después del periapsis, sin embargo, se separa hasta 68\,000~km
en 11 días y no reproduce la caída de la energía jovicéntrica real: la $\vinf$ equivalente pasa
de 7.650 a 7.578~km/s, y en el modelo solo baja a 7.640 por la marea solar
(figura~\ref{fig:voy}b).

Para encontrar la causa se calculó la aceleración real de la sonda (derivada numérica de la
velocidad del kernel) menos la gravitacional del modelo (figura~\ref{fig:voy}a). Con Júpiter
puntual aparecen los sobrevuelos de Ganímedes y Europa. Con el modelo ampliado el residuo cae a
$6\times10^{-4}$~mm/s$^2$, salvo en un intervalo de 76.6 minutos que empieza 130 minutos después
del máximo acercamiento, en el que la sonda tiene una aceleración constante de 2.83~mm/s$^2$
($\Delta v = 13.0$~m/s) casi opuesta a su velocidad respecto a Júpiter. Una aceleración así no
es gravitacional; es el empuje de una maniobra. La crónica de la misión \cite{morrison1980}
describe, en efecto, un encendido de 76 minutos unas dos horas después del máximo acercamiento.
El sobrevuelo real no fue del todo balístico.

\begin{figure}[tb]
\centering
\includegraphics[width=0.85\textwidth]{figuras/fig09_voyager2_satelites.png}
\caption{Voyager~2 cerca del máximo acercamiento a Júpiter. (a) Aceleración real menos la
gravitacional del modelo, con Júpiter puntual y con satélites y $J_2$: lo que queda es un
encendido de 76 minutos. (b) Energía jovicéntrica, expresada como
$\sqrt{2\varepsilon_{\mathrm J}}$, de los datos reales y de tres versiones del Modelo 2.}
\label{fig:voy}
\end{figure}

La caída de 0.072~km/s queda repartida así: marea solar 0.010, satélites y $J_2$ 0.029 y
maniobra 0.033. Con los tres efectos el modelo termina en 7.5780~km/s (real: 7.5777), el error
en el periapsis baja a 310~km y el tramo de 11 días posterior a la maniobra, integrado desde el
estado real, se reproduce con 6~km de error (790~km con Júpiter puntual). La masa puntual,
despreciable para el diseño preliminar ($0.4\,\%$ de $\vinf$), era lo que limitaba al Modelo 2.

\video{1JzqB87UpJlRrGls2nMdCAo0X2l7jum-V}{la trayectoria real de Voyager~2 y la maniobra}
\subsection{Uso del asistente de IA}

Se usó un asistente de IA (Claude) para escribir la mayor parte del código, borradores de
texto y la deducción de la formulación. Siguiendo la metodología de las 4D \cite{dakan2025},
cada afirmación cuantitativa se aceptó solo tras una verificación independiente; la bitácora
del repositorio registra 32 entradas y 16 errores del asistente. Tres ejemplos: el solucionador
de Lambert fallaba en transferencias de más de \ang{180} (detectado al comparar con integración
numérica); la primera comprobación de~\eqref{eq:DE} con datos reales daba un acuerdo fortuito
del $0.6\,\%$, porque omitía el potencial de Júpiter y usaba el $n$ medio, dos errores del
$5\,\%$ que se compensaban (detectado al rehacer el balance de energía); y el
asistente atribuyó la caída de energía de Voyager~2 solo a los satélites, cuando la mitad es la
maniobra (detectado con la aceleración residual y confirmado con una fuente histórica).

\section{Conclusiones y recomendaciones}\label{sec:conclusiones}

En el sistema de Júpiter se conservan, de forma aproximada, la energía y el momento angular
jovicéntricos, y con ellos $|\vec v_\infty|$; en el sistema del Sol no se conserva ni la energía
ni el momento angular de la sonda, y la ganancia sale del movimiento orbital del planeta. En el
sistema que rota con Júpiter se conserva exactamente $h = E - n\ell_z$ si la órbita es circular.
Con los datos de Voyager~2 la relación $\Delta E = n\Delta\ell_z$ se cumple al $0.4\,\%$ solo con
la velocidad angular instantánea de Júpiter, que en 1979 estaba cerca del afelio.

Las cónicas empalmadas con efemérides reales reproducen las ventanas de 1977 y predicen el
periapsis de las Voyager con error menor al $4\,\%$. Su error en la ganancia de energía es menor
al $2\,\%$ para $\vinf\gtrsim 8$~km/s y supera el $10\,\%$ para $\vinf\lesssim 3$~km/s; en ese
régimen sirven como primera aproximación y deben corregirse con integración numérica. El
Modelo~2 reproduce la trayectoria real de Voyager~2 solo si Júpiter no es una masa puntual y se
incluye la maniobra de la sonda. Falta repetir el barrido en tres dimensiones y con la órbita excéntrica de Júpiter.

\begin{thebibliography}{99}\small\setlength{\itemsep}{0pt}
\bibitem{flandro1966} G.~A. Flandro. Fast reconnaissance missions to the outer solar system
utilizing energy derived from the gravitational field of Jupiter. \emph{Astronautica Acta}
12, 329--337 (1966).
\bibitem{kohlhase1977} C.~E. Kohlhase y P.~A. Penzo. Voyager mission description.
\emph{Space Science Reviews} 21, 77--101 (1977).
\bibitem{curtis2020} H.~D. Curtis. \emph{Orbital Mechanics for Engineering Students}.
Butterworth-Heinemann, 4.\textsuperscript{a} ed., 2020.
\bibitem{goldstein2002} H.~Goldstein, C.~Poole y J.~Safko. \emph{Classical Mechanics}.
Addison-Wesley, 3.\textsuperscript{a} ed., 2002.
\bibitem{murray1999} C.~D. Murray y S.~F. Dermott. \emph{Solar System Dynamics}.
Cambridge University Press, 1999.
\bibitem{rein2012} H.~Rein y S.-F. Liu. REBOUND: an open-source multi-purpose $N$-body code
for collisional dynamics. \emph{Astronomy \& Astrophysics} 537, A128 (2012).
\bibitem{rein2015} H.~Rein y D.~S. Spiegel. IAS15: a fast, adaptive, high-order integrator
for gravitational dynamics. \emph{MNRAS} 446, 1424--1437 (2015).
\bibitem{folkner2009} W.~M. Folkner, J.~G. Williams y D.~H. Boggs. The planetary and lunar
ephemeris DE 421. \emph{IPN Progress Report} 42-178 (2009).
\bibitem{archinal2018} B.~A. Archinal \emph{et al.} Report of the IAU Working Group on Cartographic Coordinates and Rotational Elements: 2015. \emph{Celestial Mechanics and Dynamical Astronomy} 130, 22 (2018).
\bibitem{jacobson} R.~A. Jacobson. JUP230: solución de las órbitas de los satélites de Júpiter. Jet Propulsion Laboratory, NASA.
\bibitem{standish} E.~M. Standish. Keplerian elements for approximate positions of the major planets. JPL Solar System Dynamics. \url{https://ssd.jpl.nasa.gov/planets/approx_pos.html}.
\bibitem{cross} D.~Cross. Astronomy Engine (software), versión 2.1.19.
\url{https://github.com/cosinekitty/astronomy}.
\bibitem{acton1996} C.~H. Acton. Ancillary data services of NASA's Navigation and Ancillary
Information Facility. \emph{Planetary and Space Science} 44, 65--70 (1996).
\bibitem{pds} PDS Ring-Moon Systems Node. Voyager mission.
\url{https://pds-rings.seti.org/voyager/mission/}.
\bibitem{morrison1980} D.~Morrison y J.~Samz. \emph{Voyage to Jupiter}. NASA SP-439,
Washington D.C., 1980, cap.~7.
\bibitem{dakan2025} R.~Dakan y J.~Feller. The AI fluency framework: a descriptive and normative
model of human-AI interaction for higher education (2025). \url{https://aifluencyframework.org}.
\end{thebibliography}

\end{document}
